\thispagestyle{appendixheader}
\stepcounter{app}
\setcounter{app_fig}{1}
\setcounter{app_tab}{1}
\setcounter{equation}{0}
\renewcommand\theequation{附\arabic{app}-\arabic{equation}}
% \renewcommand\theequation{\Alph{app}.\arabic{equation}}

\setcounter{chapter}{0}
\setcounter{section}{0}
\thispagestyle{appendixheader}
\chapter*{附录一\quad 圣彼得堡悖论}\label{chap:app1}
\addcontentsline{toc}{chapter}{附录一\quad 圣彼得堡悖论}
圣彼得堡悖论最初由Nicolas Bernoulli于1713年提出，游戏规则如下：反复抛掷一枚公平硬币，直至首次出现正面为止；若第$n$次抛掷时首次出现正面，参与者获得$2^{n-1}$达克特。该游戏的期望货币收益为：
\begin{equation}
    \sum_{n=1}^{\infty} \frac{1}{2^n} \cdot 2^{n-1} = \frac{1}{2} + \frac{1}{2} + \cdots = \infty
\end{equation}
按照期望货币值最大化准则，理性决策者应愿意支付任意有限金额参与这一游戏。然而，现实中几乎无人愿意支付超过二十达克特。

实际支付意愿与期望值之间的悬殊差距是对当时主流决策框架的根本性挑战：若以期望货币值作为决策准则，则该框架无法解释人类最基本的经济直觉——这正是"圣彼得堡悖论"的核心张力所在。

Bernoulli提出的解法\citep{bernoulli1954exposition}是将决策目标从期望货币值替换为期望效用。引入对数效用函数$u(w) = \log(w)$后，游戏的期望效用为：
\begin{equation}
    \sum_{n=1}^{\infty} \frac{1}{2^n} \cdot \log(2^{n-1}) = \log 2 \cdot \sum_{n=1}^{\infty} \frac{n-1}{2^n} = \log 2 < \infty
\end{equation}
期望效用收敛为有限值，从而预测人们只愿支付有限金额，悖论得以消解。这一解法的关键在于对数函数的凹性：随财富增加，货币的边际效用递减，使得极低概率的巨额奖励对决策的实际驱动力远小于其名义期望值所暗示的程度。凹性效用函数由此成为描述风险厌恶的标准工具，并奠定了现代期望效用理论的基础。

\chapter*{附录二\quad VNM效用理论}\label{chap:app2}
\addcontentsline{toc}{chapter}{附录二\quad VNM效用理论}

von Neumann和Morgenstern\citep{vonneumann1953theory}将Bernoulli的效用直觉置于严格的公理体系之上：若决策者的偏好满足以下四条公理，则必然存在一个效用函数$u$，使其在不确定结果之间的选择等价于最大化期望效用$\mathbb{E}[u]$。

\begin{itemize}
    \item 完备性：对任意两个结果$A$、$B$，$A \succ B$、$B \succ A$、$A \sim B$三者之一必然成立。偏好关系在全域上有定义，不存在决策者无从比较的情形。
    \item 传递性：若$A \succ B$且$B \succ C$，则$A \succ C$。偏好关系在逻辑上保持一致，不存在循环偏好。完备性与传递性共同构成偏好的基本理性约束：前者要求决策者对任何选项都持有立场，后者要求这些立场不相互矛盾。
    \item 连续性：若$A \succ B \succ C$，则存在某个概率$p \in (0,1)$，使决策者对混合彩票$(A,\,p;\,C,\,1-p)$与确定性结果$B$无差异。这排除了偏好中存在词典序的可能性：没有任何结果好到无论概率多低都值得追求，也没有任何结果坏到无论概率多低都必须回避。
    \item 独立性：若$A \succ B$，则对任意结果$C$与概率$p \in (0,1)$，有$(A,\,p;\,C,\,1-p) \succ (B,\,p;\,C,\,1-p)$。在两个彩票中引入相同的背景成分$C$不应改变原有偏好次序——$C$对两者的影响完全对称，偏好应仅由$A$与$B$的差异决定。
\end{itemize}

VNM框架的贡献在于将"理性决策者应最大化期望效用"从经验猜想提升为公理推论：只要偏好满足上述理性约束，期望效用表示的存在性便是逻辑必然。若效用函数为凹函数，则直接蕴含边际效用递减与风险规避；Bernoulli的对数效用不过是满足凹性条件的一个具体实例。

\chapter*{附录三\quad 符号列表}\label{chap:app3}
\addcontentsline{toc}{chapter}{附录三\quad 符号列表}
\begin{table}[H]
    \centering
    \zihao{5}
    \setlength{\tabcolsep}{6pt}
    \renewcommand{\arraystretch}{1.15}
    \begin{tabularx}{\textwidth}{p{0.18\textwidth}X}
        \toprule
        \textbf{符号} & \textbf{说明} \\
        \midrule
        $w\, /\,Wealth$ & 财富水平（连续变量）\\
        $\tilde{W}_5$ & 财富水平（类别变量，5类）\\
        $\tilde{W}_{24}$ & 财富水平（类别变量，24类）\\
        $u\, /\,u(\cdot)$ & 效用 / 效用函数\\
        $ES$ & 赚取速度\\
        $e$ & 作为下标指示赚取速度\\
        $R$ & 奖励量\\
        $s$ & 作为下标指示消费对应的奖励量\\
        $SS$ & 消费速度\\
        $c$ & 选择\\
        $Choice_{ES}$ & 物品域选择（赚取/消费）\\
        $Choice_{LR}$ & 动作域选择（左/右）\\
        $ER$ & 汇率/交换率\\
        $\alpha$ & 边际效用曲率指数\\
        $\gamma$ & 选择敏感性参数\\
        $b$ & 价值偏置\\
        $V_E$ & 赚取价值（随财富变化，由效用模型估计）\\
        $V_E^{fixed}$ & 赚取价值（不随财富变化，由效用模型估计）\\
        $V_E^{Non}$ & 赚取价值（随财富变化，由非参数模型估计）\\
        $V_S$ & 消费价值（随财富变化，由效用模型估计）\\
        $V_S^{fixed}$ & 消费价值（不随财富变化，由效用模型估计）\\
        $V_S^{Non}$ & 消费价值（随财富变化，由非参数模型估计）\\
        $\Delta V$ & 价值差（赚取-消费）\\
        $V_{Chosen}$ & 被选价值\\
        \bottomrule
    \end{tabularx}
\end{table}

% \chapter{附录中的图表}{
% \stepcounter{app}
% \setcounter{app_fig}{1}
% \setcounter{app_tab}{1}


% 附表测试

% \begin{apptab}[!htbp]
%     \bicaption{\enspace 这是一个样表}{\enspace This is a sample table}
%     \stepcounter{app_tab}
%     \label{apptab:1}
%     \centering
%     \footnotesize% fontsize
%     \setlength{\tabcolsep}{4pt}% column separation
%     \renewcommand{\arraystretch}{1.2}%row space 
%     \begin{tabular}{lcccccccc}
%         \hline
%         行号 & \multicolumn{8}{c}{跨多列的标题}\\
%         %\cline{2-9}% partial hline from column i to column j
%         \hline
%         Row 1 & $1$ & $2$ & $3$ & $4$ & $5$ & $6$ & $7$ & $8$\\
%         \hline
%     \end{tabular}
% \end{apptab}


% \begin{apptab}[!htbp]
%     \bicaption{\enspace 这是一个样表}{\enspace This is a sample table}
%     \stepcounter{app_tab}
%     \label{apptab:2}
%     \centering
%     \footnotesize% fontsize
%     \setlength{\tabcolsep}{4pt}% column separation
%     \renewcommand{\arraystretch}{1.2}%row space 
%     \begin{tabular}{lcccccccc}
%         \hline
%         行号 & \multicolumn{8}{c}{跨多列的标题}\\
%         %\cline{2-9}% partial hline from column i to column j
%         \hline
%         Row 1 & $1$ & $2$ & $3$ & $4$ & $5$ & $6$ & $7$ & $8$\\
%         \hline
%     \end{tabular}
% \end{apptab}

% 附图测试

% \begin{appfig}[!htbp]
%     \centering
%     \includegraphics[width=0.40\textwidth]{c06h06}
%     \bicaption{\enspace 这是一个样图}{\enspace This is a sample figure}
%     \fignote{对图片的注释}
    
%     \label{appfig:1}
%     \stepcounter{app_fig}
% \end{appfig}

